% ============================================
% CONT-03 (EN)
% English translation — SPEC-935-R355
% ============================================
\section*{\textbf{CONT-03} --- You Smell It}
\markboth{CONT-03 --- You Smell It}{}

\noindent \textit{Sensory contamination}

\noindent You felt it.

\noindent At some point while reading, the smell arrived. You did not identify it right away. You thought it was the new book. You thought it was the room. You thought it was your imagination.

\noindent \textbf{Stop.}

\noindent Right now. Stop reading. Close your eyes. Breathe in deeply.

\noindent What did you feel?

\noindent Describe it mentally. Do not skip this step. If you skipped it, go back and do it. I know you skipped it. I will wait.

\noindent ...

\noindent Are you back?

\noindent You did not go back, did you. You kept reading. You thought the instruction was not for you. It was. Everything in this book is for you. But it is all right. Disobedience is also part of the cycle. You do not follow instructions. Joaquim did not either. That is why he ended up in the trench.

\vspace{0.5em}\noindent\rule{\textwidth}{0.5pt}\vspace{0.5em}

\noindent The smell is sweet.

\noindent It is not a clean sweetness. It is a dirty sweetness. It recalls a wilted flower inside a closed room. It recalls the bedclothes of someone who was sick for a long time. It recalls the breath of a dog that ate something rotten.

\noindent You raised the book to your nose. Just a little. Just to confirm.

\noindent Point at yourself: did you raise the book, yes or no?

\noindent If yes: the smell is there. Between the pages. In the ink. In the paper. You did not imagine it. You tasted it. The smell entered your nostrils and went down to the bottom of your lungs. In there, it settled in the alveoli. It is going to stay there. Every breath will bring back a molecule of the smell.

\noindent If no: lie. You raised it. Everyone raises it. It is a reflex. The smell is described in enough detail for your brain to simulate it. Your brain has already simulated it. The smell is in your limbic system. You do not need to feel it in your nose to feel it in your brain. You have already felt it.

\vspace{0.5em}\noindent\rule{\textwidth}{0.5pt}\vspace{0.5em}

\noindent \textbf{Instruction:} Run your index finger over the page. Lightly. Like a caress.

\noindent Did you feel the texture? The paper is smooth. Too smooth for an ordinary book. Almost damp. It is not damp --- it is \textit{warm}. The paper is at the temperature of your body. As if it were warming up with your touch.

\noindent Run your finger over it again. Now run it over your lips.

\noindent Did you do it? If you did, you felt the taste. It is metallic. Subtle. The taste of licking an envelope. But with a bitter undertone. The taste of old ink. The taste of dried blood.

\noindent If you did not, you are missing the experience. But do not worry. Later, when you are distracted, your finger will return to the page by itself. And you will lick it. Instinctively. Like someone checking.

\vspace{0.5em}\noindent\rule{\textwidth}{0.5pt}\vspace{0.5em}

\noindent You are breathing faster.

\noindent Notice. Your chest is rising and falling at a faster pace than when you started reading this fragment. Your pupils are dilated. Your heart rate has increased.

\noindent This is not suspense.

\noindent This is \textbf{contamination}.

\noindent Your body is responding to a stimulus that is not in the environment. The stimulus is in the text. The text is in your head. Your head is in your body. The body does not distinguish between real and imagined. For the body, imagining the smell is \textit{feeling} the smell. Imagining fear is \textit{having} fear.

\noindent The creature does not need to exist in the real world.
It exists in the text.
The text exists in your head.
For your body, it is real.

\vspace{0.5em}\noindent\rule{\textwidth}{0.5pt}\vspace{0.5em}

\noindent Now, lower the book.

\noindent Put it in your lap. Look at your right hand. The hand holding the book. Look at the fingers. At the joints. At the skin.

\noindent Do you see how pale the skin is? How the veins stand out more? The hand is colder than it should be.

\noindent It is not your imagination.

\noindent Peripheral circulation is shutting down. It is a stress response. Your body is preparing for a threat. What threat? The text. But the text is no longer on the page. The text is inside you. The danger is not outside. The danger is \textit{the memory of what you read}.

\noindent Your body is afraid of you.

\vspace{0.5em}\noindent\rule{\textwidth}{0.5pt}\vspace{0.5em}

\noindent Feel the weight of the book in your lap.

\noindent It is heavier than before. It is not your imagination. The weight increased. You can check if you like --- weigh the book in your hand. It is heavier.

\noindent It is not the book that weighs more.
It is your arm that weighs less.
Muscle fatigue is a symptom of prolonged sympathetic activation.
You have been on alert for too long.
Your body is tired of being afraid.

\vspace{0.5em}\noindent\rule{\textwidth}{0.5pt}\vspace{0.5em}

\noindent \textbf{Now, one last instruction.}

\noindent Look at the upper right corner of the room. Where the wall meets the ceiling.

\noindent What did you see?

\noindent Nothing? Are you sure?

\noindent Look again. Longer. Fix your eye on the dark corner for 10 seconds. Count.

\noindent 1... 2... 3... 4... 5...

\noindent Your brain is hunting for patterns where there are none. It is combing the darkness for something that might be a threat. It will not find anything. But it will \textit{almost} find something. A shadow. A movement. A gleam.

\noindent It may have been a reflection.

\noindent It may have been the creature.

\noindent You decide.

\vspace{0.5em}\noindent\rule{\textwidth}{0.5pt}\vspace{0.5em}

\noindent The smell is stronger now.

\noindent You are getting used to it. Or the smell is increasing. Hard to tell. Both things happen at the same time when contamination begins.

\noindent The Creature is closer.

\noindent Not in the room. On the page. In your head. In the space between what you read and what you imagine.

\noindent You will not get rid of the smell.

\noindent It will be on your pillow tonight.
In tomorrow's breakfast.
On the skin of whoever sleeps beside you.

\noindent You will wake in the middle of the night with your nose full of this smell.
There will be no book nearby.
There will be no explanation.
The smell will be there. Inside you.

\noindent \textbf{Because the trench is not a place.}
\textbf{The trench is a state.}
\textbf{You entered it on the first page.}
\textbf{It has no exit.}

\vspace{0.5em}\noindent\rule{\textwidth}{0.5pt}\vspace{0.5em}

\noindent \textit{--- Breathe deeply. The smell is part of you now.}

\noindent \textit{--- Run your hand over your face. Feel the temperature of your skin. It is colder.}

\noindent \textit{--- Now turn the page. The creature has waited long enough.}

\vspace{0.5em}
\noindent\textit{\textbf{↪ Routes:} \rota{CONT-04} \rotasep \rota{MEM-05}}

\vspace{0.5em}\noindent\rule{\textwidth}{0.5pt}\vspace{0.5em}

\subsection*{CONTAMINATION SCALE --- UPDATE 2}

\noindent \textit{The archive noticed your breathing.}

\begin{moltabletwo}
\hline
Points & Diagnosis \\
\hline
0 & You have not read yet. Go back now. \\
\hline
1--3 & You have read too much. The first symptoms begin in 3... 2... \\
\hline
4--6 & You are contaminated. The scale no longer measures you --- it measures \textit{what grows in you}. \\
\hline
7--10 & You already know what comes next. The mirror will show you. \\
\hline
\end{moltabletwo}

\noindent \textit{The creature is closer than you think.}

\vspace{0.5em}\noindent\rule{\textwidth}{0.5pt}\vspace{0.5em}
